\documentclass[11pt,a4paper]{article}
\usepackage[utf8]{inputenc}
\usepackage[french]{babel}
\usepackage{amsmath,amssymb,amsthm}
\usepackage{geometry}
\geometry{hmargin=2.5cm,vmargin=3cm}
\usepackage{tikz}
\usetikzlibrary{cd,positioning,shapes}
\usepackage{listings}
\usepackage{xcolor}

\lstset{
    language=Python,
    basicstyle=\ttfamily\small,
    keywordstyle=\color{blue}\bfseries,
    stringstyle=\color{red},
    commentstyle=\color{gray}\itshape,
    numbers=left,
    numberstyle=\tiny\color{gray},
    breaklines=true,
    frame=single,
    showstringspaces=false,
    literate={é}{{\'e}}1 {è}{{\`e}}1 {à}{{\`a}}1 {ê}{{\^e}}1 {î}{{\^i}}1 {ô}{{\^o}}1 {û}{{\^u}}1 {ù}{{\`u}}1 {ç}{{\c{c}}}1 {É}{{\'E}}1 {È}{{\`E}}1 {Ê}{{\^E}}1
}

\title{\Huge INTRODUCTION À LA RÉDUCTION DE JORDAN ET STRUCTURE DES NILPOTENTS \\ \large Cours magistral, exercices corrigés et travaux pratiques}
\author{\Large Charles EDOU NZE}
\date{\today}

\begin{document}
\maketitle
\tableofcontents
\newpage

\part{Cours Magistral}

\section{Introduction à la réduction de Jordan et structure des nilpotents}
\subsection{Introduction}
L'algèbre linéaire commence souvent par une promesse séduisante : celle de la diagonalisation. L'idée que l'on puisse trouver un système de coordonnées dans lequel un endomorphisme complexe agit comme une simple dilatation indépendante le long de chaque axe. Cependant, très tôt, on découvre que cette promesse est fragile. Toutes les matrices ne sont pas diagonalisables. Que se passe-t-il alors lorsque l'espace refuse de se scinder en directions purement indépendantes ? Que faire des endomorphismes qui recèlent en eux une forme d'asymétrie structurelle, une force de cisaillement que de simples vecteurs propres ne peuvent capturer ?

C'est ici qu'intervient le concept de nilpotence, puis de réduction de Jordan. Historiquement, Camille Jordan, à la fin du XIXe siècle, cherchait à comprendre la structure fine des équations différentielles linéaires. Lorsqu'un système présente des racines multiples dans son polynôme caractéristique, la solution ne s'exprime plus seulement comme une somme d'exponentielles pures, mais fait apparaître des termes de la forme $t^k e^{\lambda t}$. Cet effet de résonance est la signature spectrale d'un bloc de Jordan.

Pour comprendre la réduction de Jordan, il faut visualiser l'endomorphisme non pas comme un ensemble de ressorts indépendants, mais comme une cascade ou une chaîne de montage. Un opérateur nilpotent agit comme un broyeur : peu importe le vecteur que vous y insérez, si vous le passez assez de fois à travers la machine, il finira broyé et réduit au vecteur nul. Les vecteurs ne sont plus des entités autonomes mais forment des chaînes de dépendance, où chaque itération de l'endomorphisme pousse un vecteur vers le suivant, jusqu'à l'annihilation finale.

La forme de Jordan est ainsi le compromis ultime : c'est la matrice la plus diagonale possible. Elle révèle que tout opérateur se décompose en deux actions simultanées et qui commutent : une dilatation pure et un cisaillement nilpotent.

\subsection{Définitions, Théorèmes et Exemples}
\subsubsection{Endomorphismes Nilpotents}
Soit $E$ un $\mathbb{K}$-espace vectoriel. Un endomorphisme $u \in \mathcal{L}(E)$ est dit nilpotent s'il existe un entier $k \in \mathbb{N}^*$ tel que $u^k = 0_{\mathcal{L}(E)}$. L'entier minimal $p \in \mathbb{N}^*$ vérifiant $u^p = 0_{\mathcal{L}(E)}$ est appelé l'indice de nilpotence de $u$.

\begin{itemize}
\item $E$ : Un espace vectoriel sur le corps $\mathbb{K}$.
\item $u \in \mathcal{L}(E)$ : Un opérateur linéaire de $E$ dans $E$.
\item $k \in \mathbb{N}^*$ : Un entier naturel non nul, représentant le nombre d'itérations.
\item $u^k = u \circ u \circ \dots \circ u$ : La composition de l'endomorphisme $u$ avec lui-même $k$ fois.
\item $0_{\mathcal{L}(E)}$ : L'endomorphisme nul.
\item L'indice $p$ : Il vérifie $u^{p-1} \neq 0$ et $u^p = 0$. Il quantifie la longueur maximale de survie d'un vecteur sous l'action de $u$.
\end{itemize}

Exemple trivial : L'endomorphisme nul $u = 0_{\mathcal{L}(E)}$. Son indice de nilpotence est $p=1$.

Exemple complexe : Dans l'espace des polynômes de degré au plus $n$, $E = \mathbb{R}_n[X]$, l'opérateur de dérivation $D(P) = P'$. Puisque chaque dérivation diminue le degré d'au moins 1, la dérivée $(n+1)$-ème d'un polynôme de degré $\leq n$ est nulle. Ainsi, $D^{n+1} = 0$. L'opérateur de dérivation est nilpotent d'indice $n+1$.

Rotation de $\pi/2$ : Dans $\mathbb{R}^2$, la rotation $R$ d'angle $\pi/2$ vérifie $R^4 = \text{Id}$. Elle n'est pas nilpotente car ses itérées forment une suite périodique.

Dimension infinie : L'opérateur de décalage à droite $S(x_0, x_1, x_2, \dots) = (0, x_0, x_1, \dots)$ sur l'espace des suites n'est pas nilpotent, bien que son opérateur adjoint, le décalage à gauche $L(x_0, x_1, x_2, \dots) = (x_1, x_2, x_3, \dots)$, soit tel que $L^k$ s'annule sur le sous-espace des suites nulles à partir du rang $k$.

\begin{center}
\begin{tikzpicture}[>=stealth, auto, node distance=2cm, main node/.style={circle,draw,font=\sffamily\Large\bfseries}]
  \node[main node] (1) {$x$};
  \node[main node] (2) [right of=1] {$u(x)$};
  \node[main node] (3) [right of=2] {$u^2(x)$};
  \node (4) [right of=3] {$\dots$};
  \node[main node] (5) [right of=4] {$0$};
  \draw[->] (1) edge node {$u$} (2);
  \draw[->] (2) edge node {$u$} (3);
  \draw[->] (3) edge node {$u$} (4);
  \draw[->] (4) edge node {$u$} (5);
\end{tikzpicture}
\end{center}

\subsubsection{Caractérisation des Nilpotents}
Soit $E$ un $\mathbb{K}$-espace vectoriel de dimension finie $n$. Soit $u \in \mathcal{L}(E)$.
Les propositions suivantes sont équivalentes :
\begin{enumerate}
\item $u$ est nilpotent.
\item Le polynôme caractéristique de $u$ est $\chi_u(X) = X^n$.
\item La seule valeur propre de $u$ dans la clôture algébrique de $\mathbb{K}$ est $0$.
\end{enumerate}

La matrice $U = \begin{pmatrix} 0 & 1 & 0 \\ 0 & 0 & 1 \\ 0 & 0 & 0 \end{pmatrix}$ vérifie $\chi_U(X) = X^3$. On calcule $U^2 = \begin{pmatrix} 0 & 0 & 1 \\ 0 & 0 & 0 \\ 0 & 0 & 0 \end{pmatrix}$ et $U^3 = 0_3$. $U$ est bien nilpotente.

L'équivalence $\chi_u(X) = X^n \iff u$ nilpotent repose sur la dimension finie. En dimension infinie, on ne peut pas utiliser le polynôme caractéristique.

\subsubsection{Blocs de Jordan}
Un bloc de Jordan de taille $k \in \mathbb{N}^*$ associé à la valeur propre $\lambda \in \mathbb{K}$ est une matrice carrée $J_k(\lambda) \in \mathcal{M}_k(\mathbb{K})$ de la forme :
$$ J_k(\lambda) = \begin{pmatrix} \lambda & 1 & 0 & \dots & 0 \\ 0 & \lambda & 1 & \ddots & \vdots \\ \vdots & \ddots & \ddots & \ddots & 0 \\ \vdots & & \ddots & \lambda & 1 \\ 0 & \dots & \dots & 0 & \lambda \end{pmatrix} $$
Autrement dit, $[J_k(\lambda)]_{i,j} = \lambda$ si $i=j$, $1$ si $j = i+1$, et $0$ sinon.

$J_1(3) = \begin{pmatrix} 3 \end{pmatrix}$, $J_2(0) = \begin{pmatrix} 0 & 1 \\ 0 & 0 \end{pmatrix}$, $J_3(-2) = \begin{pmatrix} -2 & 1 & 0 \\ 0 & -2 & 1 \\ 0 & 0 & -2 \end{pmatrix}$.

\begin{center}
\begin{tikzpicture}
    \node (O) at (0,0) {};
    \node (X) at (3,0) {};
    \node (Y) at (0,3) {};
    \draw[->, thick] (O) -- (X) node[right] {$e_1$};
    \draw[->, thick] (O) -- (Y) node[above] {$e_2$};
    \draw[->, dashed, red, thick] (1, 1) -- (1, 0) node[midway, right] {$N(e_2)=e_1$};
    \draw[->, dashed, red, thick] (2, 2) -- (2, 0);
    \node at (2.5, 1) {Cisaillement nilpotent};
\end{tikzpicture}
\end{center}

\subsection{Démonstrations}
\subsubsection{Théorème de structure nilpotente et polynôme caractéristique}
Soit $E$ un $\mathbb{K}$-espace vectoriel de dimension $n \ge 1$ et $u \in \mathcal{L}(E)$.

Sens direct : Supposons $u$ nilpotent.
\begin{enumerate}
\item Par définition, il existe $k \in \mathbb{N}^*$ tel que $u^k = 0$.
\item Soit $\lambda \in \mathbb{K}$ une valeur propre de $u$ et $x \in E \setminus \{0_E\}$ un vecteur propre associé. Ainsi $u(x) = \lambda x$.
\item Montrons par récurrence sur $j$ que $u^j(x) = \lambda^j x$.
\item Appliquons ceci pour $j=k$ : $u^k(x) = \lambda^k x$.
\item Or nous savons que $u^k = 0$, donc $u^k(x) = 0_E$.
\item Ainsi, $\lambda^k x = 0_E$.
\item Comme $x \neq 0_E$, la nullité implique que $\lambda^k = 0$.
\item Dans le corps $\mathbb{K}$, cela implique nécessairement $\lambda = 0$.
\item Le polynôme caractéristique $\chi_u(X)$ de $u$ est de degré $n$.
\item La seule racine possible dans tout corps de décomposition est $0$.
\item On déduit que $\chi_u(X) = X^n$.
\end{enumerate}

Sens réciproque : Supposons $\chi_u(X) = X^n$.
\begin{enumerate}
\item D'après le théorème de Cayley-Hamilton, tout endomorphisme $u$ annule son polynôme caractéristique : $\chi_u(u) = 0_{\mathcal{L}(E)}$.
\item Remplaçons $\chi_u$ par son expression : $\chi_u(u) = u^n = 0_{\mathcal{L}(E)}$.
\item Par définition, $u$ est donc nilpotent, d'indice de nilpotence $p \le n$.
\end{enumerate}

\subsubsection{Indice de nilpotence et sous-espaces emboîtés}
Si $u$ est nilpotent d'indice $p$, alors la suite des noyaux vérifie :
$\{0_E\} = \ker(u^0) \subsetneq \ker(u^1) \subsetneq \ker(u^2) \subsetneq \dots \subsetneq \ker(u^p) = E$.

\begin{enumerate}
\item Inclusion $\ker(u^i) \subseteq \ker(u^{i+1})$ :
   Soit $x \in \ker(u^i)$. Par définition, $u^i(x) = 0_E$.
   Appliquons $u$ : $u(u^i(x)) = u(0_E) = 0_E$.
   Donc $u^{i+1}(x) = 0_E$, ce qui signifie $x \in \ker(u^{i+1})$.
\item Caractère strict des inclusions avant l'indice $p$ :
   Supposons qu'il existe un entier $k < p$ tel que $\ker(u^k) = \ker(u^{k+1})$.
   La suite des noyaux stationne à partir du rang $k$.
\item Comme $u$ est d'indice $p$, $\ker(u^p) = E$.
\item Si la suite stationnait à $k < p$, on aurait $\ker(u^p) = \ker(u^k) = E$, ce qui signifierait que $u^k = 0$.
\item Or $p$ est le plus petit entier annulant $u$, ce qui contredit $u^k=0$ pour $k < p$.
\item Donc toutes les inclusions jusqu'à $p$ sont strictes.
\end{enumerate}

\subsection{Applications}
En intelligence artificielle, la théorie des endomorphismes nilpotents et de la forme de Jordan est fondamentale pour analyser la dynamique des Réseaux de Neurones Récurrents (RNN).

Dans un RNN sans activation, l'état caché évolue selon $h_t = W h_{t-1}$.
L'analyse de la stabilité de cette dynamique sur de longues séquences dépend de $W^t$.

Si $W$ n'est pas diagonalisable et possède des blocs de Jordan associés à des valeurs propres de module 1, la matrice $J_k(1)^t$ génère des termes polynomiaux en $t$. Plus précisément, la composante sur la sur-diagonale du $j$-ème niveau croît en $\approx t^j$.

Cette croissance polynomiale d'un système qui paraissait spectralement stable peut engendrer le phénomène d'explosion du gradient ou sa version nilpotente : un bloc de Jordan pour $\lambda=0$ agit comme un filtre FIR pur, anéantissant toute information du passé après exactement $k$ pas de temps, empêchant le réseau d'apprendre des dépendances à long terme. C'est l'essence même du problème que les architectures LSTM ont été conçues pour résoudre en court-circuitant cette multiplication matricielle itérée.


\newpage
\part{Exercices d'Application et de Concours}
\subsection*{Exercice 01 : Nilpotence basique en dimension 2 \quad $\bigstar\star\star\star\star$}
\subsubsection{Énoncé}
Soit $N \in \mathcal{M}_2(\mathbb{R})$ la matrice définie par :
$$N = \begin{pmatrix} 2 & -1 \\ 4 & -2 \end{pmatrix}$$
\begin{enumerate}
\item Calculer $N^2$. Que peut-on en déduire sur l'endomorphisme canoniquement associé à $N$ ?
\item Déterminer le polynôme caractéristique $\chi_N(X)$ de $N$ et vérifier le lien avec la nilpotence.
\end{enumerate}

\subsubsection{Corrigé Rigoureux : Démonstration Complète}
\subsubsection{Calcul de $N^2$}
Posons $N = \begin{pmatrix} 2 & -1 \\ 4 & -2 \end{pmatrix}$.
Nous calculons le produit matriciel $N \times N$ :
$$N^2 = \begin{pmatrix} 2 & -1 \\ 4 & -2 \end{pmatrix} \begin{pmatrix} 2 & -1 \\ 4 & -2 \end{pmatrix}$$
Calcul des coefficients de $N^2 = (c_{i,j})$ :
\begin{itemize}
\item $c_{1,1} = 2 \times 2 + (-1) \times 4 = 4 - 4 = 0$
\item $c_{1,2} = 2 \times (-1) + (-1) \times (-2) = -2 + 2 = 0$
\item $c_{2,1} = 4 \times 2 + (-2) \times 4 = 8 - 8 = 0$
\item $c_{2,2} = 4 \times (-1) + (-2) \times (-2) = -4 + 4 = 0$
Ainsi,
$$N^2 = \begin{pmatrix} 0 & 0 \\ 0 & 0 \end{pmatrix} = 0_{\mathcal{M}_2(\mathbb{R})}$$
L'endomorphisme canoniquement associé à $N$ s'annule pour la puissance $2$. Puisque $N \neq 0$, son indice de nilpotence est exactement $p = 2$. $N$ est une matrice nilpotente.
\end{itemize}

\subsubsection{Polynôme caractéristique}
Par définition, le polynôme caractéristique de $N$ est $\chi_N(X) = \det(X I_2 - N)$.
$$\chi_N(X) = \det \begin{pmatrix} X - 2 & 1 \\ -4 & X + 2 \end{pmatrix}$$
Calculons ce déterminant :
$$\chi_N(X) = (X - 2)(X + 2) - (1 \times (-4))$$
$$\chi_N(X) = (X^2 - 4) + 4$$
$$\chi_N(X) = X^2$$
Puisque le polynôme caractéristique est $\chi_N(X) = X^2$, le théorème de caractérisation spectrale des matrices nilpotentes en dimension $n=2$ s'applique : l'unique valeur propre de $N$ est $0$. Ceci est cohérent avec le fait que $N$ est nilpotente (et, réciproquement, en dimension finie, une unique valeur propre nulle implique la nilpotence par le théorème de Cayley-Hamilton qui donne $N^2 = 0$).


\subsection*{Exercice 02 : Indice de nilpotence d'un opérateur de dérivation \quad $\bigstar\bigstar\star\star\star$}
\subsubsection{Énoncé}
Soit $E = \mathbb{R}_3[X]$ l'espace vectoriel des polynômes à coefficients réels de degré au plus 3.
Soit $u \in \mathcal{L}(E)$ l'endomorphisme de dérivation : $\forall P \in E, u(P) = P'$.
\begin{enumerate}
\item Écrire la matrice de $u$ dans la base canonique $\mathcal{B} = (1, X, X^2, X^3)$.
\item Démontrer que $u$ est nilpotent et déterminer son indice de nilpotence $p$.
\item Calculer $\ker(u^k)$ pour $k \in \{1, 2, 3, 4\}$ et vérifier les inclusions strictes.
\end{enumerate}

\subsubsection{Corrigé Rigoureux : Démonstration Complète}
\subsubsection{Matrice de $u$ dans la base canonique}
La base canonique est $\mathcal{B} = (e_0, e_1, e_2, e_3) = (1, X, X^2, X^3)$. La dimension de $E$ est $n=4$.
Calculons l'image par $u$ de chaque vecteur de base :
\begin{itemize}
\item $u(e_0) = u(1) = 0 = 0 \cdot e_0 + 0 \cdot e_1 + 0 \cdot e_2 + 0 \cdot e_3$
\item $u(e_1) = u(X) = 1 = 1 \cdot e_0 + 0 \cdot e_1 + 0 \cdot e_2 + 0 \cdot e_3$
\item $u(e_2) = u(X^2) = 2X = 0 \cdot e_0 + 2 \cdot e_1 + 0 \cdot e_2 + 0 \cdot e_3$
\item $u(e_3) = u(X^3) = 3X^2 = 0 \cdot e_0 + 0 \cdot e_1 + 3 \cdot e_2 + 0 \cdot e_3$
\end{itemize}

La matrice $M = \text{Mat}_{\mathcal{B}}(u)$ s'écrit en disposant ces coordonnées en colonnes :
$$M = \begin{pmatrix} 0 & 1 & 0 & 0 \\ 0 & 0 & 2 & 0 \\ 0 & 0 & 0 & 3 \\ 0 & 0 & 0 & 0 \end{pmatrix}$$

\subsubsection{Nilpotence et indice}
Calculons les puissances successives de $u$ (ou de $M$) :
\begin{itemize}
\item Pour tout $P = a_0 + a_1 X + a_2 X^2 + a_3 X^3$,
\item $u(P) = P' = a_1 + 2a_2 X + 3a_3 X^2$ (degré $\le 2$)
\item $u^2(P) = P'' = 2a_2 + 6a_3 X$ (degré $\le 1$)
\item $u^3(P) = P''' = 6a_3$ (degré $0$, constante)
\item $u^4(P) = P^{(4)} = 0$ (polynôme nul)
Puisque $u^4(P) = 0$ pour tout $P \in E$, on a $u^4 = 0_{\mathcal{L}(E)}$. L'endomorphisme $u$ est donc nilpotent.
L'indice de nilpotence est le plus petit entier $p$ tel que $u^p = 0$.
Ici, $u^3(X^3) = 6 \neq 0$, donc $u^3 \neq 0_{\mathcal{L}(E)}$.
Ainsi, l'indice de nilpotence de $u$ est exactement $p = 4$.
\end{itemize}

\subsubsection{Étude des noyaux itérés}
Nous devons trouver l'ensemble des polynômes $P$ tels que $u^k(P) = 0$.
\begin{itemize}
\item $\ker(u^1) = \{ P \in E \mid P' = 0 \}$. Ce sont les polynômes constants. $\ker(u^1) = \text{Vect}(1) = \mathbb{R}_0[X]$.
\item $\ker(u^2) = \{ P \in E \mid P'' = 0 \}$. En intégrant deux fois, on trouve les polynômes de degré $\le 1$. $\ker(u^2) = \text{Vect}(1, X) = \mathbb{R}_1[X]$.
\item $\ker(u^3) = \{ P \in E \mid P''' = 0 \}$. Les polynômes de degré $\le 2$. $\ker(u^3) = \text{Vect}(1, X, X^2) = \mathbb{R}_2[X]$.
\item $\ker(u^4) = \{ P \in E \mid P^{(4)} = 0 \}$. L'espace $E$ tout entier. $\ker(u^4) = \mathbb{R}_3[X] = E$.
On vérifie la suite strictement croissante d'inclusions (jusqu'à $k=p=4$) :
$$\{0\} \subsetneq \text{Vect}(1) \subsetneq \text{Vect}(1, X) \subsetneq \text{Vect}(1, X, X^2) \subsetneq \mathbb{R}_3[X]$$
Les dimensions croissent strictement de 1 à chaque étape (1, 2, 3, 4).
\end{itemize}


\subsection*{Exercice 03 : Matrices triangulaires strictes \quad $\bigstar\bigstar\star\star\star$}
\subsubsection{Énoncé}
Soit $T_n(\mathbb{K})$ l'ensemble des matrices triangulaires supérieures strictes de taille $n \times n$ (matrices dont tous les coefficients diagonaux et sous-diagonaux sont nuls).
Démontrer rigoureusement par récurrence que le produit de $k$ matrices triangulaires strictes est une matrice dont les $k$ premières sur-diagonales sont nulles. En déduire que toute matrice triangulaire stricte de taille $n$ est nilpotente d'indice au plus $n$.

\subsubsection{Corrigé Rigoureux : Démonstration Complète}
\subsubsection{Structure du produit de matrices triangulaires strictes}
Soit $A = (a_{i,j})_{1 \le i,j \le n}$ une matrice triangulaire supérieure stricte.
Par définition, $a_{i,j} = 0$ si $j \le i$. (En d'autres termes, $a_{i,j} \neq 0 \implies j > i \implies j \ge i+1$).
Considérons la propriété $\mathcal{P}(k)$ : "Le produit de $k$ matrices triangulaires supérieures strictes $A^{(1)} A^{(2)} \dots A^{(k)}$ est une matrice $M = (m_{i,j})$ telle que $m_{i,j} = 0$ pour $j < i + k$."

\textbf{Initialisation ($k=1$) :}
Pour une seule matrice $M = A^{(1)}$, l'hypothèse est que $M$ est triangulaire stricte. Donc $m_{i,j} = 0$ pour $j \le i$, soit $j < i + 1$. La propriété $\mathcal{P}(1)$ est vraie.

\textbf{Hérédité :}
Supposons $\mathcal{P}(k)$ vraie pour un entier $k \ge 1$.
Soit $M = A^{(1)} \dots A^{(k)}$, avec $m_{i,j} = 0$ si $j < i + k$.
Soit $B = A^{(k+1)}$ une matrice triangulaire stricte, avec $b_{i,j} = 0$ si $j \le i$.
Calculons le produit $C = M B$. Ses coefficients sont $c_{i,j} = \sum_{\ell=1}^n m_{i,\ell} b_{\ell,j}$.
Nous voulons montrer que $c_{i,j} = 0$ si $j < i + k + 1$.
Analysons les termes de la somme : $m_{i,\ell} b_{\ell,j}$.
Pour qu'un terme soit non nul, il faut simultanément :
\begin{enumerate}
\item $m_{i,\ell} \neq 0 \implies \ell \ge i + k$ (d'après $\mathcal{P}(k)$).
\item $b_{\ell,j} \neq 0 \implies j \ge \ell + 1$ (car $B$ est triangulaire stricte).
En combinant ces deux inégalités strictes :
$j \ge \ell + 1 \ge (i + k) + 1 = i + k + 1$.
Ainsi, si $j < i + k + 1$, alors pour tout $\ell \in \{1, \dots, n\}$, soit $m_{i,\ell} = 0$, soit $b_{\ell,j} = 0$.
Par conséquent, la somme est nulle : $c_{i,j} = 0$.
La propriété $\mathcal{P}(k+1)$ est donc vérifiée.
Par principe de récurrence, $\mathcal{P}(k)$ est vraie pour tout $k \ge 1$.
\end{enumerate}

\subsubsection{Nilpotence}
Appliquons $\mathcal{P}(n)$ en choisissant $A^{(1)} = A^{(2)} = \dots = A^{(n)} = A$.
La matrice $M = A^n$ vérifie $m_{i,j} = 0$ si $j < i + n$.
Or, les indices des lignes $i$ et des colonnes $j$ varient entre $1$ et $n$.
La plus grande valeur possible pour $j$ est $n$, et la plus petite valeur pour $i$ est $1$.
On a toujours $j \le n < n + 1 \le i + n$.
Donc la condition $j < i + n$ est vérifiée par tous les couples $(i,j)$ de la matrice.
Par conséquent, tous les coefficients de $A^n$ sont nuls : $A^n = 0_{\mathcal{M}_n(\mathbb{K})}$.
Toute matrice triangulaire stricte de taille $n$ est donc nilpotente, d'indice de nilpotence au plus $n$.


\subsection*{Exercice 04 : Bloc de Jordan fondamental \quad $\bigstar\bigstar\bigstar\star\star$}
\subsubsection{Énoncé}
Soit $J \in \mathcal{M}_3(\mathbb{R})$ le bloc de Jordan nilpotent canonique :
$$J = \begin{pmatrix} 0 & 1 & 0 \\ 0 & 0 & 1 \\ 0 & 0 & 0 \end{pmatrix}$$
Soit $N \in \mathcal{M}_3(\mathbb{R})$ une matrice telle que $N^3 = 0$ et $N^2 \neq 0$.
\begin{enumerate}
\item Calculer les puissances successives de $J$.
\item Montrer que $\ker(N)$ est de dimension 1.
\item Démontrer qu'il existe un vecteur $x \in \mathbb{R}^3$ tel que la famille $\mathcal{B} = (N^2 x, Nx, x)$ est une base de $\mathbb{R}^3$.
\item Quelle est la matrice de l'endomorphisme associé à $N$ dans cette base $\mathcal{B}$ ? Conclure.
\end{enumerate}

\subsubsection{Corrigé Rigoureux : Démonstration Complète}
\subsubsection{Puissances de $J$}
$$J = \begin{pmatrix} 0 & 1 & 0 \\ 0 & 0 & 1 \\ 0 & 0 & 0 \end{pmatrix}$$
On calcule $J^2$ :
$$J^2 = \begin{pmatrix} 0 & 1 & 0 \\ 0 & 0 & 1 \\ 0 & 0 & 0 \end{pmatrix} \begin{pmatrix} 0 & 1 & 0 \\ 0 & 0 & 1 \\ 0 & 0 & 0 \end{pmatrix} = \begin{pmatrix} 0 & 0 & 1 \\ 0 & 0 & 0 \\ 0 & 0 & 0 \end{pmatrix}$$
On calcule $J^3$ :
$$J^3 = J^2 J = \begin{pmatrix} 0 & 0 & 1 \\ 0 & 0 & 0 \\ 0 & 0 & 0 \end{pmatrix} \begin{pmatrix} 0 & 1 & 0 \\ 0 & 0 & 1 \\ 0 & 0 & 0 \end{pmatrix} = \begin{pmatrix} 0 & 0 & 0 \\ 0 & 0 & 0 \\ 0 & 0 & 0 \end{pmatrix}$$
L'indice de nilpotence de $J$ est exactement 3.

\subsubsection{Dimension du noyau}
Soit $u$ l'endomorphisme canoniquement associé à $N$. $u$ est nilpotent d'indice 3.
Nous avons la suite stricte des noyaux :
$\{0\} \subsetneq \ker(u) \subsetneq \ker(u^2) \subsetneq \ker(u^3) = \mathbb{R}^3$
Soit $d_k = \dim(\ker(u^k))$. La suite $(d_k)$ est strictement croissante.
De plus, $d_0 = 0$ et $d_3 = 3$.
Puisque les inclusions sont strictes, les dimensions sautent d'au moins 1 à chaque étape.
Il faut faire 3 sauts pour passer de 0 à 3, donc $d_k = k$ pour tout $k \in \{0, 1, 2, 3\}$.
En particulier, $\dim(\ker(u)) = 1$.

\subsubsection{Construction de la base}
Puisque $u^2 \neq 0$, il existe au moins un vecteur $x \in \mathbb{R}^3$ tel que $u^2(x) \neq 0$.
Considérons la famille $\mathcal{B} = (u^2(x), u(x), x)$.
Montrons qu'elle est libre.
Soit $a, b, c \in \mathbb{R}$ tels que $a u^2(x) + b u(x) + c x = 0$.
Appliquons $u^2$ à cette équation :
$a u^4(x) + b u^3(x) + c u^2(x) = 0$.
Puisque $u^3 = 0$ (donc $u^4 = 0$), il reste $c u^2(x) = 0$.
Comme $u^2(x) \neq 0$, nous déduisons $c = 0$.
L'équation initiale se réduit à $a u^2(x) + b u(x) = 0$.
Appliquons $u$ :
$a u^3(x) + b u^2(x) = 0$.
De nouveau, $u^3 = 0$, donc $b u^2(x) = 0$, ce qui implique $b = 0$.
L'équation se réduit à $a u^2(x) = 0$, donc $a = 0$.
La famille $\mathcal{B}$ est libre.
Puisque son cardinal est 3 et $\dim(\mathbb{R}^3) = 3$, c'est une base de $\mathbb{R}^3$.

\subsubsection{Matrice dans la base $\mathcal{B}$}
Posons $e_1 = u^2(x)$, $e_2 = u(x)$, $e_3 = x$.
Calculons l'image par $u$ des vecteurs de cette base :
\begin{itemize}
\item $u(e_1) = u(u^2(x)) = u^3(x) = 0 = 0 \cdot e_1 + 0 \cdot e_2 + 0 \cdot e_3$
\item $u(e_2) = u(u(x)) = u^2(x) = e_1 = 1 \cdot e_1 + 0 \cdot e_2 + 0 \cdot e_3$
\item $u(e_3) = u(x) = e_2 = 0 \cdot e_1 + 1 \cdot e_2 + 0 \cdot e_3$
La matrice de $u$ dans la base $\mathcal{B}$ est donc :
$$M = \begin{pmatrix} 0 & 1 & 0 \\ 0 & 0 & 1 \\ 0 & 0 & 0 \end{pmatrix} = J$$
\textbf{Conclusion :} Tout endomorphisme nilpotent d'indice 3 en dimension 3 est semblable au bloc de Jordan canonique $J_3(0)$. Ce processus (choisir un vecteur qui ne s'annule qu'à la puissance maximale et construire la famille de ses images) s'appelle la construction d'une chaîne de Jordan.
\end{itemize}


\subsection*{Exercice 05 : Nilpotence et trace \quad $\bigstar\bigstar\bigstar\star\star$}
\subsubsection{Énoncé}
Soit $E$ un $\mathbb{K}$-espace vectoriel de dimension $n$. Soit $u \in \mathcal{L}(E)$.
On suppose que $\mathbb{K}$ est un corps de caractéristique nulle (ex: $\mathbb{R}$ ou $\mathbb{C}$).
\begin{enumerate}
\item Si $u$ est nilpotent, montrer que pour tout entier $k \ge 1$, $\text{Tr}(u^k) = 0$.
\item (Réciproque partielle) On suppose que pour tout $k \in \{1, 2, \dots, n\}$, $\text{Tr}(u^k) = 0$. En utilisant les identités de Newton (admises), montrer que $u$ est nilpotent.
\end{enumerate}

\subsubsection{Corrigé Rigoureux : Démonstration Complète}
\subsubsection{Condition nécessaire}
Si $u$ est nilpotent, alors son polynôme caractéristique est $\chi_u(X) = X^n$ (Théorème vu en cours).
Cela signifie que dans la clôture algébrique de $\mathbb{K}$ (ou dans $\mathbb{C}$ si on travaille sur $\mathbb{R}$), toutes les valeurs propres de $u$ sont nulles.
$u$ est trigonalisable dans $\mathbb{C}$, donc il existe une base dans laquelle sa matrice $T$ est triangulaire supérieure avec des zéros sur la diagonale.
Pour tout entier $k \ge 1$, la matrice $T^k$ est également triangulaire supérieure, et ses éléments diagonaux sont les éléments diagonaux de $T$ élevés à la puissance $k$, donc des zéros.
La trace étant invariante par changement de base et égale à la somme des éléments diagonaux de $T^k$, on a $\text{Tr}(u^k) = 0$ pour tout $k \ge 1$.

\subsubsection{Condition suffisante via Newton}
Soient $\lambda_1, \lambda_2, \dots, \lambda_n$ les racines complexes (valeurs propres) du polynôme caractéristique de $u$, répétées avec leur multiplicité.
On sait que $\text{Tr}(u^k) = \sum_{i=1}^n \lambda_i^k$.
Par hypothèse, $S_k = \sum_{i=1}^n \lambda_i^k = 0$ pour tout $k \in \{1, \dots, n\}$.
Considérons le polynôme caractéristique factorisé :
$\chi_u(X) = X^n - e_1 X^{n-1} + e_2 X^{n-2} - \dots + (-1)^n e_n$
où les $e_k$ sont les fonctions symétriques élémentaires des racines.
Les identités de Newton relient les sommes de puissances $S_k$ aux $e_k$ par la relation récursive (pour $k \le n$) :
$k e_k = \sum_{i=1}^k (-1)^{i-1} e_{k-i} S_i \quad (\text{avec } e_0 = 1)$
Appliquons ceci itérativement :
\begin{itemize}
\item Pour $k=1$ : $1 e_1 = S_1$. Puisque $S_1 = 0$, on a $e_1 = 0$.
\item Pour $k=2$ : $2 e_2 = e_1 S_1 - S_2 = 0 \times 0 - 0 = 0$. Comme la caractéristique du corps est nulle, $2 \neq 0$, donc $e_2 = 0$.
\item Par récurrence forte, si $e_1 = e_2 = \dots = e_{k-1} = 0$,
  alors $k e_k = S_k - e_1 S_{k-1} + \dots = 0$.
  La division par $k$ (possible en caractéristique nulle) donne $e_k = 0$.
Ainsi, pour tout $k \in \{1, \dots, n\}$, $e_k = 0$.
Le polynôme caractéristique de $u$ est donc simplement :
$\chi_u(X) = X^n$
D'après le théorème de Cayley-Hamilton, $u^n = \chi_u(u) = 0_{\mathcal{L}(E)}$.
L'endomorphisme $u$ est donc nilpotent.
\end{itemize}


\subsection*{Exercice 06 : Commutant d'un bloc de Jordan \quad $\bigstar\bigstar\bigstar\star\star$}
\subsubsection{Énoncé}
Soit $J_n \in \mathcal{M}_n(\mathbb{R})$ le bloc de Jordan nilpotent canonique (des 1 sur la sur-diagonale, 0 ailleurs).
Soit $C(J_n) = \{ M \in \mathcal{M}_n(\mathbb{R}) \mid M J_n = J_n M \}$ le commutant de $J_n$.
\begin{enumerate}
\item Démontrer que si $M$ commute avec $J_n$, alors $M$ est un polynôme en $J_n$. (Indice : évaluer les relations de commutation sur les coefficients de la matrice).
\item En déduire la dimension de l'espace vectoriel $C(J_n)$.
\end{enumerate}

\subsubsection{Corrigé Rigoureux : Démonstration Complète}
\subsubsection{Analyse de la commutation}
Soit $M = (m_{i,j})_{1 \le i,j \le n}$ une matrice.
Le produit $J_n M$ est la matrice dont la $i$-ème ligne est la $(i+1)$-ème ligne de $M$ pour $i<n$, et la dernière ligne est nulle.
$(J_n M)_{i,j} = m_{i+1, j} \quad (\text{avec } m_{n+1, j} = 0)$
Le produit $M J_n$ est la matrice dont la $j$-ème colonne est la $(j-1)$-ème colonne de $M$ pour $j>1$, et la première colonne est nulle.
$(M J_n)_{i,j} = m_{i, j-1} \quad (\text{avec } m_{i, 0} = 0)$

L'égalité $M J_n = J_n M$ s'écrit composante par composante :
Pour tout $1 \le i \le n$ et $1 \le j \le n$, $m_{i, j-1} = m_{i+1, j}$.
\begin{itemize}
\item Pour $j=1$ : $m_{i+1, 1} = m_{i, 0} = 0$. Donc la première colonne de $M$ (sauf $m_{1,1}$) est nulle : $m_{2,1} = m_{3,1} = \dots = m_{n,1} = 0$.
\item Pour $i=n$ : $m_{n, j-1} = m_{n+1, j} = 0$. Donc la dernière ligne de $M$ (sauf $m_{n,n}$) est nulle : $m_{n,1} = m_{n,2} = \dots = m_{n,n-1} = 0$.
\item Par la relation $m_{i, j-1} = m_{i+1, j}$, la matrice $M$ est constante sur chaque diagonale parallèle à la diagonale principale. Une telle matrice s'appelle une matrice de Toeplitz.
Les coefficients sous-diagonaux sont nuls (car égaux aux éléments de la première colonne qui sont nuls).
Il ne reste que des valeurs arbitraires sur la diagonale principale ($a_0$) et les sur-diagonales ($a_1, a_2, \dots, a_{n-1}$).
$M$ a donc la forme :
$$M = \begin{pmatrix} a_0 & a_1 & a_2 & \dots & a_{n-1} \\ 0 & a_0 & a_1 & \dots & a_{n-2} \\ 0 & 0 & a_0 & \ddots & \vdots \\ \vdots & \vdots & \ddots & \ddots & a_1 \\ 0 & 0 & \dots & 0 & a_0 \end{pmatrix}$$
Or, on vérifie facilement que la matrice avec des 1 sur la $k$-ème sur-diagonale est exactement $(J_n)^k$.
Par convention, $(J_n)^0 = I_n$.
Ainsi, $M$ s'écrit de manière unique comme combinaison linéaire des puissances de $J_n$ :
$M = a_0 I_n + a_1 J_n + a_2 (J_n)^2 + \dots + a_{n-1} (J_n)^{n-1}$.
Cela montre que $M$ est un polynôme en $J_n$ : $M \in \mathbb{R}[J_n]$.
\end{itemize}

\subsubsection{Dimension du commutant}
Nous avons montré que $C(J_n) \subseteq \text{Vect}(I_n, J_n, (J_n)^2, \dots, (J_n)^{n-1})$.
Réciproquement, tout polynôme en $J_n$ commute évidemment avec $J_n$ (car les puissances commutent entre elles).
Donc $C(J_n) = \text{Vect}(I_n, J_n, (J_n)^2, \dots, (J_n)^{n-1})$.
Il reste à montrer que la famille $\mathcal{F} = (I_n, J_n, \dots, (J_n)^{n-1})$ est libre.
Soit une combinaison linéaire nulle : $\sum_{k=0}^{n-1} \lambda_k (J_n)^k = 0$.
La matrice résultante est exactement la matrice de Toeplitz avec les $\lambda_k$ sur les diagonales. Pour qu'elle soit nulle, il faut que tous ses coefficients soient nuls, c'est-à-dire $\lambda_0 = \lambda_1 = \dots = \lambda_{n-1} = 0$.
La famille est donc une base de $C(J_n)$.
La dimension du commutant est donc le nombre d'éléments de cette base, soit $n$.


\subsection*{Exercice 07 : Inverse d'un bloc diagonalisable + nilpotent \quad $\bigstar\bigstar\bigstar\bigstar\star$}
\subsubsection{Énoncé}
Soit $A \in \mathcal{M}_n(\mathbb{R})$ inversible.
Supposons que $A = \lambda I_n - N$, où $\lambda \in \mathbb{R}^*$ et $N$ est une matrice nilpotente d'indice $p$.
\begin{enumerate}
\item Rappeler le développement en série entière formelle de $(1 - x)^{-1}$.
\item Utiliser ce développement pour exprimer $A^{-1}$ comme un polynôme en $N$.
\item Application : Soit le bloc de Jordan inversible $J_{\lambda} = \lambda I_3 + J$ (où $J$ est le bloc nilpotent usuel de taille 3, $\lambda \neq 0$). Calculer l'inverse exact de $J_{\lambda}$.
\end{enumerate}

\subsubsection{Corrigé Rigoureux : Démonstration Complète}
\subsubsection{Développement formel}
Pour un réel $x$ tel que $|x| < 1$, on a le développement géométrique classique :
$\frac{1}{1 - x} = \sum_{k=0}^{+\infty} x^k = 1 + x + x^2 + x^3 + \dots$

\subsubsection{Expression de $A^{-1}$}
Factorisons $A$ :
$A = \lambda I_n - N = \lambda (I_n - \frac{1}{\lambda} N)$.
Nous cherchons un inverse $B$ tel que $A B = I_n$, soit $\lambda (I_n - \frac{1}{\lambda} N) B = I_n$.
Cela s'écrit : $(I_n - \frac{1}{\lambda} N) B = \frac{1}{\lambda} I_n$.
Par analogie avec le développement scalaire, posons formellement :
$S = \sum_{k=0}^{+\infty} (\frac{1}{\lambda} N)^k$
Puisque $N$ est nilpotente d'indice $p$, on a $N^k = 0$ pour tout $k \ge p$.
La somme infinie est donc en réalité une somme finie :
$S = \sum_{k=0}^{p-1} \frac{1}{\lambda^k} N^k$
Vérifions que $S$ est l'inverse cherché en calculant $(I_n - \frac{1}{\lambda} N) S$ :
$(I_n - \frac{1}{\lambda} N) \sum_{k=0}^{p-1} (\frac{1}{\lambda} N)^k = \sum_{k=0}^{p-1} (\frac{1}{\lambda} N)^k - \frac{1}{\lambda} N \sum_{k=0}^{p-1} (\frac{1}{\lambda} N)^k$
$= \sum_{k=0}^{p-1} (\frac{1}{\lambda} N)^k - \sum_{k=0}^{p-1} (\frac{1}{\lambda} N)^{k+1}$
Par télescopage, presque tous les termes s'annulent. Il reste :
$= (\frac{1}{\lambda} N)^0 - (\frac{1}{\lambda} N)^p = I_n - \frac{1}{\lambda^p} N^p$
Or, par définition de l'indice de nilpotence, $N^p = 0$.
Donc le produit vaut bien $I_n$.
Ainsi, $(I_n - \frac{1}{\lambda} N)^{-1} = \sum_{k=0}^{p-1} \frac{1}{\lambda^k} N^k$.
On en déduit que $A^{-1} = \frac{1}{\lambda} (I_n - \frac{1}{\lambda} N)^{-1} = \sum_{k=0}^{p-1} \frac{1}{\lambda^{k+1}} N^k$.
L'inverse $A^{-1}$ est bien un polynôme en $N$.

\subsubsection{Application au bloc de Jordan}
Soit $J_{\lambda} = \lambda I_3 + J$, avec $J = \begin{pmatrix} 0 & 1 & 0 \\ 0 & 0 & 1 \\ 0 & 0 & 0 \end{pmatrix}$.
Ici $A = \lambda I_3 - (-J)$. Donc on applique la formule avec $N = -J$ et $p=3$.
$A^{-1} = \frac{1}{\lambda} I_3 + \frac{1}{\lambda^2} (-J) + \frac{1}{\lambda^3} (-J)^2$.
Calculons les termes :
\begin{itemize}
\item $J = \begin{pmatrix} 0 & 1 & 0 \\ 0 & 0 & 1 \\ 0 & 0 & 0 \end{pmatrix}$
\item $J^2 = \begin{pmatrix} 0 & 0 & 1 \\ 0 & 0 & 0 \\ 0 & 0 & 0 \end{pmatrix}$
Ainsi, l'inverse est :
$$J_{\lambda}^{-1} = \begin{pmatrix} \frac{1}{\lambda} & -\frac{1}{\lambda^2} & \frac{1}{\lambda^3} \\ 0 & \frac{1}{\lambda} & -\frac{1}{\lambda^2} \\ 0 & 0 & \frac{1}{\lambda} \end{pmatrix}$$
Ce résultat illustre parfaitement comment inverser explicitement des blocs diagonaux non diagonalisables.
\end{itemize}


\subsection*{Exercice 08 : Décomposition d'un noyau itéré et stabilité \quad $\bigstar\bigstar\bigstar\bigstar\star$}
\subsubsection{Énoncé}
Soit $E$ un espace vectoriel de dimension finie, et $u \in \mathcal{L}(E)$ nilpotent d'indice $p$.
Pour $1 \le k \le p$, on définit les sous-espaces $N_k = \ker(u^k)$ et $I_k = \text{Im}(u^k)$.
\begin{enumerate}
\item Montrer que $E = N_p$.
\item Démontrer que $u(N_{k}) \subseteq N_{k-1}$ pour tout $1 \le k \le p$.
\item Démontrer que si $F$ est un sous-espace vectoriel de $E$ supplémentaire de $N_{p-1}$ dans $N_p$ (soit $N_p = N_{p-1} \oplus F$), alors l'application restreinte $u : F \to u(F)$ est un isomorphisme.
\item En déduire que $u(F) \cap N_{p-2} = \{0_E\}$.
\end{enumerate}

\subsubsection{Corrigé Rigoureux : Démonstration Complète}
\subsubsection{$E = N_p$}
Par définition de l'indice de nilpotence, $p$ est l'entier tel que $u^p = 0_{\mathcal{L}(E)}$.
Pour tout $x \in E$, on a $u^p(x) = 0_E$.
Donc par définition du noyau, $x \in \ker(u^p) = N_p$. Ainsi $E \subseteq N_p$.
Comme $N_p \subseteq E$ trivialement, $E = N_p$.

\subsubsection{Stabilité descendante}
Soit $k \in \{1, \dots, p\}$. Soit $y \in u(N_k)$.
Par définition de l'image directe, il existe $x \in N_k$ tel que $y = u(x)$.
Comme $x \in N_k$, on a $u^k(x) = 0_E$.
Calculons l'action de $u^{k-1}$ sur $y$ :
$u^{k-1}(y) = u^{k-1}(u(x)) = u^k(x) = 0_E$.
Donc par définition, $y \in \ker(u^{k-1}) = N_{k-1}$.
L'inclusion $u(N_k) \subseteq N_{k-1}$ est démontrée.

\subsubsection{Isomorphisme de restriction}
Soit $F$ tel que $N_p = N_{p-1} \oplus F$.
L'application $u_F : F \to u(F)$ définie par $x \mapsto u(x)$ est trivialement surjective par définition de son espace d'arrivée.
Montrons qu'elle est injective. Soit $x \in F$ tel que $u_F(x) = 0_E$.
Alors $u(x) = 0_E$, ce qui implique $x \in \ker(u) = N_1$.
Comme $p \ge 1$, et sachant que la suite des noyaux est croissante, on a $N_1 \subseteq N_{p-1}$.
Donc $x \in N_{p-1}$.
Or $x$ appartient aussi à $F$. Donc $x \in N_{p-1} \cap F$.
Par définition de la somme directe $N_{p-1} \oplus F$, l'intersection est réduite à $\{0_E\}$.
Donc $x = 0_E$.
L'application $u_F$ est linéaire, surjective et injective, c'est donc un isomorphisme d'espaces vectoriels. On en déduit notamment que $\dim(u(F)) = \dim(F)$.

\subsubsection{Intersection avec $N_{p-2}$}
Soit $y \in u(F) \cap N_{p-2}$.
Puisque $y \in u(F)$, il existe un unique $x \in F$ tel que $y = u(x)$.
Puisque $y \in N_{p-2}$, on a $u^{p-2}(y) = 0_E$.
Remplaçons $y$ : $u^{p-2}(u(x)) = 0_E$, c'est-à-dire $u^{p-1}(x) = 0_E$.
Cela signifie que $x \in \ker(u^{p-1}) = N_{p-1}$.
Or, nous savons que $x \in F$.
Donc $x \in N_{p-1} \cap F = \{0_E\}$ par la somme directe.
Ainsi, $x = 0_E$, ce qui implique $y = u(0_E) = 0_E$.
L'intersection est donc bien réduite au vecteur nul : $u(F) \cap N_{p-2} = \{0_E\}$.
\textit{(Ce résultat est fondamental pour construire la base de Jordan "du haut vers le bas").}


\subsection*{Exercice 09 : Matrices semblables et invariants de similitude \quad $\bigstar\bigstar\bigstar\bigstar\star$}
\subsubsection{Énoncé}
Soient les deux matrices nilpotentes suivantes dans $\mathcal{M}_3(\mathbb{R})$ :
$$A = \begin{pmatrix} 0 & 1 & 1 \\ 0 & 0 & 1 \\ 0 & 0 & 0 \end{pmatrix}, \quad B = \begin{pmatrix} 0 & 1 & 0 \\ 0 & 0 & 1 \\ 0 & 0 & 0 \end{pmatrix}$$
\begin{enumerate}
\item Calculer les polynômes caractéristiques $\chi_A$ et $\chi_B$. Sont-ils égaux ?
\item Calculer $A^2$ et $A^3$. Déterminer l'indice de nilpotence de $A$.
\item $A$ et $B$ sont-elles semblables ? Justifier de manière formelle.
\end{enumerate}

\subsubsection{Corrigé Rigoureux : Démonstration Complète}
\subsubsection{Polynômes caractéristiques}
Les matrices $A$ et $B$ sont toutes deux des matrices triangulaires supérieures strictes.
Les éléments diagonaux sont donc les valeurs propres, et elles sont toutes nulles.
Le déterminant de $X I_3 - M$ pour une matrice triangulaire supérieure stricte $M$ est simplement le produit des termes de la diagonale $(X - 0)$.
Donc $\chi_A(X) = X^3$ et $\chi_B(X) = X^3$.
Les polynômes caractéristiques sont identiques.

\subsubsection{Puissances et indice de nilpotence de $A$}
Calculons $A^2$ :
$$A^2 = \begin{pmatrix} 0 & 1 & 1 \\ 0 & 0 & 1 \\ 0 & 0 & 0 \end{pmatrix} \begin{pmatrix} 0 & 1 & 1 \\ 0 & 0 & 1 \\ 0 & 0 & 0 \end{pmatrix}$$
\begin{itemize}
\item Ligne 1, Col 3 : $0\times 1 + 1\times 1 + 1\times 0 = 1$.
Les autres coefficients sont nuls par les propriétés des matrices triangulaires strictes (Exercice 3).
$$A^2 = \begin{pmatrix} 0 & 0 & 1 \\ 0 & 0 & 0 \\ 0 & 0 & 0 \end{pmatrix}$$
$A^2 \neq 0$.
Par le théorème de Cayley-Hamilton, ou par calcul direct de $A^3 = A^2 A$, on sait que $A^3 = 0_3$.
L'indice de nilpotence de $A$ est donc $3$.
\end{itemize}

\subsubsection{Similitude}
Deux matrices sont semblables si et seulement si elles représentent le même endomorphisme dans des bases différentes. Si elles sont semblables, elles partagent les mêmes invariants de similitude : trace, déterminant, polynôme caractéristique, polynôme minimal, et indice de nilpotence.

La matrice $B$ est le bloc de Jordan canonique $J_3(0)$. Nous avons démontré dans l'exercice 4 que son indice de nilpotence est 3.
L'indice de nilpotence de $A$ est également 3.
Soit $u$ l'endomorphisme de $\mathbb{R}^3$ dont la matrice dans la base canonique est $A$.
$u$ est nilpotent d'indice $p=3$.
Comme démontré dans l'Exercice 4, tout endomorphisme nilpotent de dimension 3 et d'indice 3 admet un vecteur $x$ tel que $(u^2(x), u(x), x)$ est une base. Dans cette base, la matrice de $u$ est exactement $B = J_3(0)$.
Par conséquent, il existe une matrice de passage $P$ inversible telle que $A = P B P^{-1}$.
Les matrices $A$ et $B$ sont donc rigoureusement semblables.


\subsection*{Exercice 10 : Construction explicite de la base de Jordan \quad $\bigstar\bigstar\bigstar\bigstar\bigstar$}
\subsubsection{Énoncé}
Considérons l'endomorphisme $u$ de $\mathbb{R}^3$ dont la matrice dans la base canonique est :
$$A = \begin{pmatrix} 2 & -1 & 2 \\ 5 & -3 & 3 \\ -1 & 0 & -2 \end{pmatrix}$$
\begin{enumerate}
\item Calculer le polynôme caractéristique de $A$ et vérifier qu'il s'écrit $\chi_A(X) = (X+1)^3$.
\item En déduire que la matrice $N = A + I_3$ est nilpotente.
\item Déterminer l'indice de nilpotence $p$ de $N$ et la suite des noyaux $\ker(N^k)$.
\item Construire explicitement une base de Jordan $\mathcal{B}$ pour $N$, et donner la matrice de passage $P$.
\end{enumerate}

\subsubsection{Corrigé Rigoureux : Démonstration Complète}
\subsubsection{Polynôme caractéristique}
Calculons $\chi_A(X) = \det(X I_3 - A)$.
$$\chi_A(X) = \det \begin{pmatrix} X - 2 & 1 & -2 \\ -5 & X + 3 & -3 \\ 1 & 0 & X + 2 \end{pmatrix}$$
Développons par rapport à la dernière ligne :
$\chi_A(X) = 1 \cdot \det \begin{pmatrix} 1 & -2 \\ X+3 & -3 \end{pmatrix} - 0 + (X+2) \cdot \det \begin{pmatrix} X-2 & 1 \\ -5 & X+3 \end{pmatrix}$
$= (-3 - (-2)(X+3)) + (X+2)((X-2)(X+3) - (-5))$
$= (-3 + 2X + 6) + (X+2)(X^2 + X - 6 + 5)$
$= (2X + 3) + (X+2)(X^2 + X - 1)$
$= 2X + 3 + (X^3 + X^2 - X + 2X^2 + 2X - 2)$
$= X^3 + 3X^2 + 3X + 1$
On reconnaît l'identité remarquable du cube : $\chi_A(X) = (X + 1)^3$.

\subsubsection{Nilpotence de $N$}
Posons $N = A - (-1)I_3 = A + I_3$.
L'unique valeur propre de $A$ est $-1$.
La matrice $N$ a pour polynôme caractéristique $\chi_N(X) = \chi_A(X - 1) = ((X-1)+1)^3 = X^3$.
D'après le théorème de Cayley-Hamilton, $N^3 = 0$. $N$ est donc nilpotente.

\subsubsection{Indice de nilpotence et noyaux}
Calculons $N$ :
$$N = A + I_3 = \begin{pmatrix} 3 & -1 & 2 \\ 5 & -2 & 3 \\ -1 & 0 & -1 \end{pmatrix}$$
Calculons $N^2$ :
$$N^2 = \begin{pmatrix} 3 & -1 & 2 \\ 5 & -2 & 3 \\ -1 & 0 & -1 \end{pmatrix} \begin{pmatrix} 3 & -1 & 2 \\ 5 & -2 & 3 \\ -1 & 0 & -1 \end{pmatrix} = \begin{pmatrix} 2 & -1 & 1 \\ 2 & -1 & 1 \\ -2 & 1 & -1 \end{pmatrix}$$
On voit que $N^2 \neq 0_3$. Comme on sait que $N^3 = 0_3$, l'indice de nilpotence de $N$ est $p=3$.
Les dimensions des noyaux successifs sont $d_1 = 1$, $d_2 = 2$, $d_3 = 3$ (car la suite est strictement croissante).

\subsubsection{Construction de la base de Jordan}
L'indice de nilpotence étant $p=3$ en dimension 3, il n'y a qu'un seul bloc de Jordan $J_3(0)$ pour $N$.
Nous devons trouver un vecteur $x$ tel que $N^2 x \neq 0$.
Au vu de $N^2$, choisissons le vecteur de la base canonique $e_1 = \begin{pmatrix} 1 \\ 0 \\ 0 \end{pmatrix}$.
On a $N^2 e_1 = \text{première colonne de } N^2 = \begin{pmatrix} 2 \\ 2 \\ -2 \end{pmatrix}$. Il est bien non nul.
Construisons la chaîne de Jordan ascendante (pour avoir des 1 sur la sur-diagonale) :
\begin{itemize}
\item $v_1 = N^2 e_1 = \begin{pmatrix} 2 \\ 2 \\ -2 \end{pmatrix}$
\item $v_2 = N e_1 = \text{première colonne de } N = \begin{pmatrix} 3 \\ 5 \\ -1 \end{pmatrix}$
\item $v_3 = e_1 = \begin{pmatrix} 1 \\ 0 \\ 0 \end{pmatrix}$
\end{itemize}

Vérifions l'action de $N$ sur cette base $\mathcal{B} = (v_1, v_2, v_3)$ :
\begin{itemize}
\item $N(v_1) = N^3 e_1 = 0 = 0 v_1 + 0 v_2 + 0 v_3$
\item $N(v_2) = N^2 e_1 = v_1 = 1 v_1 + 0 v_2 + 0 v_3$
\item $N(v_3) = N e_1 = v_2 = 0 v_1 + 1 v_2 + 0 v_3$
\end{itemize}

Dans cette base, la matrice de $N$ est exactement le bloc de Jordan $J_3(0)$.
La matrice de passage de la base canonique vers $\mathcal{B}$ est la concaténation en colonnes :
$$P = \begin{pmatrix} 2 & 3 & 1 \\ 2 & 5 & 0 \\ -2 & -1 & 0 \end{pmatrix}$$
On conclut que $P^{-1} N P = J_3(0)$.
Par ailleurs, comme $A = -I_3 + N$, on a $P^{-1} A P = P^{-1} (-I_3 + N) P = -I_3 + P^{-1} N P = -I_3 + J_3(0)$.
La matrice $A$ est donc réductible à la forme de Jordan $J = \begin{pmatrix} -1 & 1 & 0 \\ 0 & -1 & 1 \\ 0 & 0 & -1 \end{pmatrix}$. $\blacksquare$




\newpage
\part{Travaux Pratiques et Simulations Algorithmiques}
\subsection{TP 01 : Opérations matricielles basiques et test de nilpotence}
\subsubsection{Objectif}
Implémenter la multiplication matricielle purement en Python, sans aucune bibliothèque externe.
Utiliser cette fonction pour écrire un algorithme capable de déterminer si une matrice est nilpotente et, le cas échéant, de trouver son indice de nilpotence.

\subsubsection{Implémentation et Validation Mathématique}
\begin{lstlisting}
# --- IMPLÉMENTATION PURE PYTHON ---

def matrix_multiply(A, B):
    """
    Multiplie deux matrices A et B représentées par des listes de listes.
    Vérifie les dimensions.
    """
    rows_A = len(A)
    cols_A = len(A[0])
    rows_B = len(B)
    cols_B = len(B[0])

    if cols_A != rows_B:
        raise ValueError("Dimensions incompatibles pour la multiplication.")

    # Initialisation de la matrice résultat avec des zéros
    C = [[0.0 for _ in range(cols_B)] for _ in range(rows_A)]

    for i in range(rows_A):
        for j in range(cols_B):
            for k in range(cols_A):
                C[i][j] += A[i][k] * B[k][j]

    return C

def is_zero_matrix(A, tol=1e-9):
    """Vérifie si tous les coefficients d'une matrice sont nuls à une tolérance près."""
    for row in A:
        for val in row:
            if abs(val) > tol:
                return False
    return True

def find_nilpotence_index(A, max_iter=100):
    """
    Détermine l'indice de nilpotence d'une matrice A.
    Retourne l'indice, ou None si la matrice n'est pas nilpotente après max_iter itérations.
    Pour une matrice n x n, si elle est nilpotente, son indice est au plus n.
    """
    n = len(A)
    if is_zero_matrix(A):
        return 1

    current_power = A
    for k in range(2, min(n + 1, max_iter + 1)):
        current_power = matrix_multiply(current_power, A)
        if is_zero_matrix(current_power):
            return k

    return None

# --- VALIDATION ET ASSERTIONS MATHÉMATIQUES ---

if __name__ == "__main__":
    # Test 1 : Matrice nilpotente d'indice 3
    # M = [[0, 1, 0], [0, 0, 1], [0, 0, 0]]
    M1 = [
        [0.0, 1.0, 0.0],
        [0.0, 0.0, 1.0],
        [0.0, 0.0, 0.0]
    ]
    idx1 = find_nilpotence_index(M1)
    assert idx1 == 3, f"Erreur: indice attendu 3, obtenu {idx1}"
    print("Test 1 passé : Bloc de Jordan de taille 3 correctement identifié, indice = 3.")

    # Test 2 : Matrice nilpotente de l'exercice 1
    # N = [[2, -1], [4, -2]]
    M2 = [
        [2.0, -1.0],
        [4.0, -2.0]
    ]
    idx2 = find_nilpotence_index(M2)
    assert idx2 == 2, f"Erreur: indice attendu 2, obtenu {idx2}"
    print("Test 2 passé : Matrice de l'exo 1 nilpotente, indice = 2.")

    # Test 3 : Matrice non nilpotente (Identité)
    I3 = [
        [1.0, 0.0, 0.0],
        [0.0, 1.0, 0.0],
        [0.0, 0.0, 1.0]
    ]
    idx3 = find_nilpotence_index(I3)
    assert idx3 is None, f"Erreur: I3 ne devrait pas être nilpotente."
    print("Test 3 passé : Identité non nilpotente identifiée.")
\end{lstlisting}


\subsection{TP 02 : Trace itérée et condition de nilpotence}
\subsubsection{Objectif}
L'exercice 5 montre que sur un corps de caractéristique nulle, $u$ est nilpotent si et seulement si $\text{Tr}(u^k) = 0$ pour tout $k \in \{1, \dots, n\}$.
Implémenter le calcul de la trace et vérifier cette équivalence algorithmiquement sur des matrices données.

\subsubsection{Implémentation et Validation Mathématique}
\begin{lstlisting}
# --- IMPLÉMENTATION PURE PYTHON ---

def matrix_multiply(A, B):
    n = len(A)
    C = [[0.0] * n for _ in range(n)]
    for i in range(n):
        for j in range(n):
            C[i][j] = sum(A[i][k] * B[k][j] for k in range(n))
    return C

def trace(A):
    """Calcule la trace d'une matrice carrée."""
    n = len(A)
    return sum(A[i][i] for i in range(n))

def check_nilpotence_via_trace(A, tol=1e-9):
    """
    Vérifie si Tr(A^k) = 0 pour k allant de 1 à n.
    C'est une condition nécessaire et suffisante de nilpotence en car 0.
    """
    n = len(A)
    current_power = A
    for k in range(1, n + 1):
        tr = trace(current_power)
        if abs(tr) > tol:
            return False # Non nilpotent
        if k < n:
            current_power = matrix_multiply(current_power, A)
    return True

# --- VALIDATION ET ASSERTIONS MATHÉMATIQUES ---

if __name__ == "__main__":
    # Test 1 : Matrice nilpotente complexe
    # M = [[2, -1, 2], [5, -3, 3], [-1, 0, -2]] est N = A + I de l'exo 10
    # Posons N = [[3, -1, 2], [5, -2, 3], [-1, 0, -1]]
    N = [
        [3.0, -1.0, 2.0],
        [5.0, -2.0, 3.0],
        [-1.0, 0.0, -1.0]
    ]
    is_nilp = check_nilpotence_via_trace(N)
    assert is_nilp == True, "Erreur: N devrait être identifiée comme nilpotente par la méthode des traces."
    print("Test 1 passé : Matrice nilpotente validée via critère des traces.")

    # Test 2 : Matrice de rotation (non nilpotente)
    # Cos(pi/2) = 0, Sin(pi/2) = 1
    R = [
        [0.0, -1.0],
        [1.0,  0.0]
    ]
    # Tr(R) = 0, mais Tr(R^2) = -2 != 0
    is_nilp2 = check_nilpotence_via_trace(R)
    assert is_nilp2 == False, "Erreur: La rotation n'est pas nilpotente."
    print("Test 2 passé : Matrice de rotation rejetée (Tr(R^2) != 0).")
\end{lstlisting}


\subsection{TP 03 : Calcul du commutant d'une matrice}
\subsubsection{Objectif}
L'exercice 6 étudie le commutant d'un bloc de Jordan. Nous allons implémenter un algorithme qui, pour une matrice $A$ donnée, construit un système d'équations linéaires pour trouver toutes les matrices $M$ telles que $A M - M A = 0$.
\textit{(On se limitera à construire la matrice du système linéaire $L(M) = AM - MA$ vue comme un vecteur aplati).}

\subsubsection{Implémentation et Validation Mathématique}
\begin{lstlisting}
# --- IMPLÉMENTATION PURE PYTHON ---

def kronecker_product_system(A):
    """
    Construit la matrice de l'opérateur L(M) = AM - MA dans la base canonique aplatie.
    Si M est aplatie par lignes, le terme (AM - MA)_{i,j} est :
    sum_k (A_{i,k} M_{k,j}) - sum_k (M_{i,k} A_{k,j})
    """
    n = len(A)
    # L'opérateur L est une matrice de taille n^2 x n^2
    L = [[0.0] * (n * n) for _ in range(n * n)]

    for i in range(n):
        for j in range(n):
            row_idx = i * n + j

            # AM : sum over k of A_{i,k} M_{k,j}
            # M_{k,j} correspond à l'index aplati col_idx_1 = k * n + j
            for k in range(n):
                col_idx_1 = k * n + j
                L[row_idx][col_idx_1] += A[i][k]

            # -MA : -sum over k of M_{i,k} A_{k,j}
            # M_{i,k} correspond à l'index aplati col_idx_2 = i * n + k
            for k in range(n):
                col_idx_2 = i * n + k
                L[row_idx][col_idx_2] -= A[k][j]

    return L

def trace_matrix(L):
    """Affiche la matrice pour l'inspection."""
    for row in L:
        print([round(v, 2) for v in row])

# --- VALIDATION ET ASSERTIONS MATHÉMATIQUES ---

if __name__ == "__main__":
    # Test avec J2(0)
    J2 = [
        [0.0, 1.0],
        [0.0, 0.0]
    ]

    L = kronecker_product_system(J2)
    # Pour J2, n=2, L est 4x4.
    # L(M) = J2 M - M J2.
    # Si M = [[a, b], [c, d]], L(M) = [[c, d-a], [0, -c]]
    # Ordre de M aplati : a, b, c, d (indices 0, 1, 2, 3)
    # L(M)[0] = c  => ligne 0 = [0, 0, 1, 0]
    # L(M)[1] = d-a => ligne 1 = [-1, 0, 0, 1]
    # L(M)[2] = 0  => ligne 2 = [0, 0, 0, 0]
    # L(M)[3] = -c => ligne 3 = [0, 0, -1, 0]

    expected_L = [
        [ 0.0,  0.0,  1.0,  0.0],
        [-1.0,  0.0,  0.0,  1.0],
        [ 0.0,  0.0,  0.0,  0.0],
        [ 0.0,  0.0, -1.0,  0.0]
    ]

    assert L == expected_L, "Erreur dans la construction du système de commutation pour J2."
    print("Test 1 passé : La matrice du système linéaire de commutation (Kronecker) est exacte.")
\end{lstlisting}


\subsection{TP 04 : Algorithme de recherche d'une chaîne de Jordan (Dimension 3)}
\subsubsection{Objectif}
Implémenter l'algorithme construit dans l'Exercice 10 : étant donné une matrice nilpotente d'indice $p=3$ en dimension 3, trouver un vecteur $x$ tel que $N^2 x \neq 0$, puis générer la base $\mathcal{B} = (N^2 x, Nx, x)$.
Construire la matrice de passage $P$ et valider que $P^{-1} N P = J_3(0)$.

\subsubsection{Implémentation et Validation Mathématique}
\begin{lstlisting}
# --- IMPLÉMENTATION PURE PYTHON ---

def mat_vec_mult(A, x):
    """Multiplie une matrice A par un vecteur colonne x."""
    n = len(A)
    y = [0.0] * n
    for i in range(n):
        y[i] = sum(A[i][j] * x[j] for j in range(n))
    return y

def matrix_multiply(A, B):
    n = len(A)
    C = [[0.0] * n for _ in range(n)]
    for i in range(n):
        for j in range(n):
            C[i][j] = sum(A[i][k] * B[k][j] for k in range(n))
    return C

def determinant_3x3(P):
    """Déterminant explicite 3x3 (Règle de Sarrus)."""
    return (P[0][0]*P[1][1]*P[2][2] + P[0][1]*P[1][2]*P[2][0] + P[0][2]*P[1][0]*P[2][1]
          - P[0][2]*P[1][1]*P[2][0] - P[0][1]*P[1][0]*P[2][2] - P[0][0]*P[1][2]*P[2][1])

def invert_3x3(P):
    """Inverse une matrice 3x3 par la comatrice."""
    det = determinant_3x3(P)
    if abs(det) < 1e-9:
        raise ValueError("Matrice non inversible")

    C = [[0.0]*3 for _ in range(3)]
    C[0][0] = P[1][1]*P[2][2] - P[1][2]*P[2][1]
    C[0][1] = -(P[1][0]*P[2][2] - P[1][2]*P[2][0])
    C[0][2] = P[1][0]*P[2][1] - P[1][1]*P[2][0]

    C[1][0] = -(P[0][1]*P[2][2] - P[0][2]*P[2][1])
    C[1][1] = P[0][0]*P[2][2] - P[0][2]*P[2][0]
    C[1][2] = -(P[0][0]*P[2][1] - P[0][1]*P[2][0])

    C[2][0] = P[0][1]*P[1][2] - P[0][2]*P[1][1]
    C[2][1] = -(P[0][0]*P[1][2] - P[0][2]*P[1][0])
    C[2][2] = P[0][0]*P[1][1] - P[0][1]*P[1][0]

    # Transposée de la comatrice / det
    Inv = [[0.0]*3 for _ in range(3)]
    for i in range(3):
        for j in range(3):
            Inv[i][j] = C[j][i] / det
    return Inv

def build_jordan_basis(N):
    """
    Trouve une chaîne de Jordan pour N nilpotente d'indice 3 dans R^3.
    """
    N2 = matrix_multiply(N, N)

    # Testons les vecteurs de la base canonique
    basis_candidates = [
        [1.0, 0.0, 0.0],
        [0.0, 1.0, 0.0],
        [0.0, 0.0, 1.0]
    ]

    x = None
    for cand in basis_candidates:
        N2_x = mat_vec_mult(N2, cand)
        # Si N^2 x n'est pas le vecteur nul
        if max(abs(val) for val in N2_x) > 1e-9:
            x = cand
            break

    if x is None:
        raise ValueError("L'indice de nilpotence n'est pas 3.")

    v1 = mat_vec_mult(N2, x)
    v2 = mat_vec_mult(N, x)
    v3 = x

    # P est la concaténation en colonnes de v1, v2, v3
    P = [
        [v1[0], v2[0], v3[0]],
        [v1[1], v2[1], v3[1]],
        [v1[2], v2[2], v3[2]]
    ]
    return P

# --- VALIDATION ET ASSERTIONS MATHÉMATIQUES ---

if __name__ == "__main__":
    # Matrice N de l'exercice 10
    N = [
        [3.0, -1.0, 2.0],
        [5.0, -2.0, 3.0],
        [-1.0, 0.0, -1.0]
    ]

    P = build_jordan_basis(N)
    P_inv = invert_3x3(P)

    # Calcul de P^{-1} N P
    Temp = matrix_multiply(N, P)
    Jordan_form = matrix_multiply(P_inv, Temp)

    # Vérifions que c'est J3(0)
    expected_J3 = [
        [0.0, 1.0, 0.0],
        [0.0, 0.0, 1.0],
        [0.0, 0.0, 0.0]
    ]

    for i in range(3):
        for j in range(3):
            assert abs(Jordan_form[i][j] - expected_J3[i][j]) < 1e-9, f"Erreur coef ({i},{j})"

    print("Test passé : Chaîne de Jordan construite, changement de base validé mathématiquement (P^{-1}NP = J3).")
\end{lstlisting}


\subsection{TP 05 : Calcul de l'inverse via série formelle de Neumann}
\subsubsection{Objectif}
L'exercice 7 établit que si $A = \lambda I_n - N$ avec $N$ nilpotente d'indice $p$, alors $A^{-1} = \sum_{k=0}^{p-1} \frac{1}{\lambda^{k+1}} N^k$.
Écrire un algorithme qui inverse un tel bloc de Jordan de manière exacte (sans pivot de Gauss) en utilisant cette série finie.

\subsubsection{Implémentation et Validation Mathématique}
\begin{lstlisting}
# --- IMPLÉMENTATION PURE PYTHON ---

def matrix_add(A, B):
    n = len(A)
    return [[A[i][j] + B[i][j] for j in range(n)] for i in range(n)]

def matrix_scale(c, A):
    n = len(A)
    return [[c * A[i][j] for j in range(n)] for i in range(n)]

def matrix_multiply(A, B):
    n = len(A)
    C = [[0.0] * n for _ in range(n)]
    for i in range(n):
        for j in range(n):
            C[i][j] = sum(A[i][k] * B[k][j] for k in range(n))
    return C

def get_identity(n):
    return [[1.0 if i == j else 0.0 for j in range(n)] for i in range(n)]

def invert_jordan_block_neumann(lmbda, N, p):
    """
    Calcule (lambda * I - N)^{-1} = sum_{k=0}^{p-1} (1 / lambda^{k+1}) N^k
    Ici, le bloc de Jordan classique est de la forme J = lambda * I + J_nilp.
    Donc on applique avec N = -J_nilp.
    """
    n = len(N)
    inverse = [[0.0]*n for _ in range(n)]

    current_power = get_identity(n) # N^0

    for k in range(p):
        coef = 1.0 / (lmbda ** (k + 1))
        term = matrix_scale(coef, current_power)
        inverse = matrix_add(inverse, term)

        # Prépare la puissance suivante
        current_power = matrix_multiply(current_power, N)

    return inverse

# --- VALIDATION ET ASSERTIONS MATHÉMATIQUES ---

if __name__ == "__main__":
    # Test avec J_lambda = lambda I + J_nilp pour n=3
    # Posons lambda = 2.0
    lmbda = 2.0
    # J_nilp = [[0, 1, 0], [0, 0, 1], [0, 0, 0]]
    J_nilp = [
        [0.0, 1.0, 0.0],
        [0.0, 0.0, 1.0],
        [0.0, 0.0, 0.0]
    ]
    # Nous devons inverser lmbda * I_3 - (-J_nilp)
    N = matrix_scale(-1.0, J_nilp) # N = [[0, -1, 0], [0, 0, -1], [0, 0, 0]]

    # Calcul par Neumann, p=3 (indice de nilpotence de J_nilp)
    Inv = invert_jordan_block_neumann(lmbda, N, 3)

    # La matrice A à inverser est:
    # A = [[2, 1, 0], [0, 2, 1], [0, 0, 2]]
    A = [
        [2.0, 1.0, 0.0],
        [0.0, 2.0, 1.0],
        [0.0, 0.0, 2.0]
    ]

    # Vérification: A * Inv doit être l'identité
    Prod = matrix_multiply(A, Inv)
    I3 = get_identity(3)

    for i in range(3):
        for j in range(3):
            assert abs(Prod[i][j] - I3[i][j]) < 1e-9, f"Erreur d'inversion coef ({i},{j})"

    print("Test passé : Inversion d'un bloc de Jordan non-diagonalisable par série de Neumann exacte (sans Gauss).")
\end{lstlisting}




\end{document}
